% year: תש"פ
% type: מבחן
% semester: א
% moed: א
% part: ב
% number: 1
% points: 9
% categories: ivt, mvt, continuity, multiple-choice
% source: exams/מבחנים/תשפ/מועד א תשפ.pdf

\begin{question}
תהי $f(x)$ רציפה בקטע $[a,b]$. מה מהבאים חייב להיות נכון?
\begin{enumerate}
  \item[(א)] $f(x)$ לא גזירה בקטע $(a,b)$
  \item[(ב)] $f(x)$ מקיימת את התנאי המשפט לגרנז' בקטע $[a,b]$
  \item[(ג)] $f(x)$ מקיימת את התנאי המשפט רול בקטע $[a,b]$
  \item[(ד)] $f(x)$ גזירה בקטע $(a,b)$
  \item[(ה)] קיים $M$ כך ש-$f(x)<M$ לכל $x$ ששייך ל-$[a,b]$
\end{enumerate}
\end{question}

\begin{hint}
איזה משפט מבטיח משהו לכל פונקציה רציפה בקטע סגור?
\end{hint}

\begin{solution}
\textbf{התשובה הנכונה: (ה).}

לפי \textbf{משפט ויירשטראס הראשון}, פונקציה רציפה בקטע סגור $[a,b]$ חסומה בו. לכן קיים $K$ כך ש-$f(x)\le K$ לכל $x\in[a,b]$, ואז $M=K+1$ מקיים $f(x)<M$ לכל $x\in[a,b]$.

שאר התשובות אינן חייבות להתקיים:
\begin{itemize}
  \item (א) לא נכון: למשל $f(x)=x$ רציפה וגזירה ב-$(a,b)$.
  \item (ב) ו-(ד) לא נכונות: משפט לגרנז' דורש גם גזירות ב-$(a,b)$, ופונקציה רציפה אינה חייבת להיות גזירה. למשל בקטע $[-1,1]$ הפונקציה $f(x)=|x|$ רציפה אך אינה גזירה ב-$0$.
  \item (ג) לא נכון: משפט רול דורש גם גזירות וגם $f(a)=f(b)$; למשל $f(x)=x$ רציפה אך $f(a)\ne f(b)$.
\end{itemize}
\end{solution}
